% !TeX program = xelatex
% !TeX encoding = UTF-8
\documentclass[UTF8]{ctexart}
\usepackage[a4paper, margin=2cm]{geometry}
\usepackage{amsmath,amssymb}
\usepackage{xcolor}
\usepackage{hyperref}

\ctexset{
  section = {
    format = \Large\bfseries\raggedright
  }
}

\hypersetup{
  colorlinks=true,
  urlcolor=blue
}

\title{Syllabus for Quantum Mechanics}
\author{Dongxue Qu}
\date{Fall 2026}

\begin{document}
\maketitle

\section*{Course Information}
\begin{itemize}
  \item Course: Quantum Mechanics
  \item Course level: Undergraduate
  \item Semester: Fall 2026
  \item Course website: \href{https://qudx54632.github.io/qudx.github.io/teaching/qm/}{https://qudx54632.github.io/qudx.github.io/teaching/qm/}
  \item Instructor: Dongxue Qu (曲东雪)
  \item Email: \href{mailto:dqu@cdut.edu.cn}{\texttt{dqu@cdut.edu.cn}}
  \item Teaching assistant: Min Sun (孙敏)
  \item Teaching assistant email: \href{mailto:3257887316@qq.com}{\texttt{3257887316@qq.com}}
  \item Office hours: Wednesdays, 16:00--17:00
  \item Location: 北翼楼，5教8楼 5824（芙蓉餐厅对面）
\end{itemize}

Emails will normally be answered during regular working hours.

\section*{Course Meeting Times}
\begin{itemize}
  \item Total number of lectures: 24.
  \item Class times: Mondays and Wednesdays, Periods 3--4 (1.5 hours per lecture).
  \item The first lecture will take place on September 2. The detailed lecture schedule will be updated on the course website throughout the semester.
\end{itemize}

\section*{Prerequisites}
Students should be familiar with calculus, linear algebra, differential equations, and classical mechanics.

\section*{Description}
This course introduces the foundations and essential techniques of quantum mechanics. Topics include quantum states and observables, wave functions and the Schr\"odinger equation, operators and quantum measurement, one-dimensional quantum systems, angular momentum and spin, central potentials, approximation methods, and identical particles.

\section*{Recommended Textbooks}
\begin{itemize}
  \item Griffiths, David J., and Darrell F. Schroeter. \textit{Introduction to Quantum Mechanics}. 3rd ed., Cambridge University Press, 2018.
  \item 曾谨言：《量子力学教程》，科学出版社。
\end{itemize}

Additional lecture notes, assignments, and resources will be posted on the course website during the semester.

\section*{Topics}
The course provides a standard undergraduate-level introduction to quantum mechanics. The main subject areas are:
\begin{itemize}
  \item Physical motivation and fundamental principles of quantum theory
  \item Quantum states, observables, measurement, and the uncertainty principle
  \item Quantum dynamics and stationary-state problems
  \item Representative one-dimensional systems, including confinement, tunneling, and scattering
  \item The harmonic oscillator and algebraic methods
  \item Rotational symmetry, angular momentum, and quantum systems with central potentials
  \item Many-particle systems, spin, periodic structures, and nonclassical correlations
\end{itemize}

\section*{Learning Objectives}
By the end of the course, students will be able to:
\begin{itemize}
  \item understand the postulates and mathematical framework of quantum mechanics;
  \item solve the Schr\"odinger equation for standard quantum systems;
  \item work with operators, commutators, expectation values, and quantum measurements;
  \item analyze angular momentum and spin; and
  \item apply approximation methods to systems without exact solutions.
\end{itemize}

\section*{Grading}
The course grade will be based on the following components:
\begin{itemize}
  \item Homework: 25\%
  \item Class participation: 10\%
  \item Final exam: 60\%
  \item Class performance (attendance and engagement): 5\%
\end{itemize}

\section*{Final Exam}
The date, time, format, and location of the final examination will be announced during the semester.

\section*{Homework}
The homework policies are as follows:
\begin{itemize}
  \item Assignments will be posted on the course website. The due date for each assignment will be provided when the assignment is released.
  \item All solutions must be handwritten and clearly legible.
  \item Late homework will receive half credit. For example, a score of 90\% on a late submission will be recorded as 45\%.
  \item The average homework score will be used to determine the homework component of the final grade.
\end{itemize}

\noindent
\textbf{Homework is an important part of this course.} Students are encouraged to work through the problems independently. They may discuss the problems with classmates or consult external resources; however, all submitted solutions must reflect each student's own work.

\noindent
\textbf{Plagiarism is a serious offense.} Do not copy another person's work or submit work that is not your own. If plagiarism is detected, the grade for that homework will be zero.

\end{document}
